% !TeX program = xelatex
% !TeX program = xelatex
% ===========================================================================
%  tech_common.tex -- 两份技术文档共用的导言区
% ===========================================================================
\documentclass[11pt,a4paper]{ctexart}

\usepackage[top=2.4cm,bottom=2.4cm,left=2.4cm,right=2.4cm]{geometry}
\usepackage{amsmath,amssymb,amsthm}
\usepackage{graphicx}
\usepackage{booktabs}
\usepackage{longtable}
\usepackage{array}
\usepackage{enumitem}
\usepackage{xcolor}
\usepackage{listings}
\usepackage{tcolorbox}
\usepackage{caption}
\usepackage{fancyhdr}
\usepackage{titlesec}
\usepackage{hyperref}

\hypersetup{colorlinks=true,linkcolor=blue!55!black,citecolor=blue!55!black,
            urlcolor=blue!55!black,bookmarksnumbered=true}

% ---------------------------------------------------------------- 字体/版式
\setlength{\emergencystretch}{2.5em}
\setlength{\parskip}{0.35em}
\setlength{\parindent}{2em}
\linespread{1.12}

% ---------------------------------------------------------------- 代码样式
\definecolor{codebg}{RGB}{247,247,250}
\definecolor{codekw}{RGB}{0,90,160}
\definecolor{codecm}{RGB}{0,128,96}
\definecolor{codest}{RGB}{170,60,60}
\lstset{
  basicstyle=\ttfamily\footnotesize,
  backgroundcolor=\color{codebg},
  frame=single, framerule=0.4pt, rulecolor=\color{gray!45},
  breaklines=true, breakatwhitespace=false,
  numbers=left, numberstyle=\tiny\color{gray!80}, numbersep=6pt,
  keywordstyle=\color{codekw}\bfseries, commentstyle=\color{codecm}\itshape,
  stringstyle=\color{codest}, showstringspaces=false, columns=flexible,
  tabsize=2, extendedchars=true, captionpos=b,
  xleftmargin=1.6em, framexleftmargin=1.4em
}
\renewcommand{\lstlistingname}{代码}

% ---------------------------------------------------------------- 提示框
\newtcolorbox{keybox}[1][]{colback=blue!4,colframe=blue!45!black,
  boxrule=0.6pt,arc=2pt,left=5pt,right=5pt,top=4pt,bottom=4pt,
  fonttitle=\bfseries,#1}
\newtcolorbox{warnbox}[1][]{colback=orange!6,colframe=orange!55!black,
  boxrule=0.6pt,arc=2pt,left=5pt,right=5pt,top=4pt,bottom=4pt,
  fonttitle=\bfseries,#1}

% ---------------------------------------------------------------- 定理环境
\newtheorem{theorem}{定理}
\newtheorem{proposition}{命题}
\newtheorem{lemma}{引理}
\renewcommand{\proofname}{\heiti 证明}

% ---------------------------------------------------------------- 页眉页脚
\pagestyle{fancy}
\fancyhf{}
\fancyhead[L]{\small\color{gray!70}\leftmark}
\fancyhead[R]{\small\color{gray!70}\thepage}
\renewcommand{\headrulewidth}{0.3pt}
\renewcommand{\headrule}{\hbox to\headwidth{\color{gray!40}\leaders\hrule height \headrulewidth\hfill}}

\titleformat{\section}{\normalfont\heiti\Large}{\thesection}{0.6em}{}
\titleformat{\subsection}{\normalfont\heiti\large}{\thesubsection}{0.6em}{}
\titleformat{\subsubsection}{\normalfont\heiti\normalsize}{\thesubsubsection}{0.6em}{}

\captionsetup{font=small,labelfont=bf,labelsep=quad}


\title{\heiti 模拟器验证与开源对照报告\\[4pt]
       \large 64 项严格检查、3 个真实缺陷、1 个把两轮调参全部作废的陷阱\\[2pt]
       \normalsize —— 以及同类开源实现的横向调研}
\author{2026 高教社杯全国大学生数学建模竞赛 B 题}
\date{\today}

\begin{document}
\maketitle
\thispagestyle{fancy}

\begin{abstract}
\noindent
本文档回答三个问题：\textbf{我的本地模拟器到底可不可信}、\textbf{有没有同类开源实现可以对照}、
\textbf{这一轮"优化"里哪些结论是真的}。

我们写了一个 64 项检查的验证套件（\texttt{code/verify\_simulator.py}），覆盖 HTTP 协议、
虚拟时间记账、物理判据、响应契约、与开源参考的成本模型交叉对照、守恒律、增量认证的等价性。
\textbf{结果是 64/64 通过，但过程中抓出并修复了 3 个真实缺陷}：
\begin{enumerate}[leftmargin=1.8em]
  \item \texttt{robot\_id} 在服务端未指定时形同虚设（HTTP 层只检查启动参数，不检查首次 \texttt{/enter} 绑定的身份）；
  \item 示向度四舍五入到两位小数后可以突破 $\pm1^{\circ}$ 的硬上界（最大超出 $0.0024^{\circ}$）；
  \item \texttt{certification\_scan} 的"未认证位置"列表被一个早退优化\textbf{截断}，规划器长期在残缺信息上决策。
\end{enumerate}

我们还找到了一份\textbf{专门针对本题的开源工程}（官方模拟器的驱动 $+$ 离线 mock），
逐项对照了 12 个物理/时间常量与完整的成本模型：\textbf{逐步完全一致}；
唯一分歧在示向度误差模型——该 mock 每次测量重新随机，违反附件 1"同一地点重复测量结果相同"的题设，
我们的实现按题设做成了地点函数。

最后如实记录一个\textbf{把两轮参数寻优全部作废的陷阱}：算例池只生成一次并复用，
而 \texttt{Arena} 会就地清除 \texttt{Source} 对象，导致第二次以后的每次评估都跑在"已经被清空的算例"上，
于是出现了一个并不存在的"44\% 提升"。加入可复现性守卫后重跑，真实收益是 4\%（问题三）与 20\%（问题四）。
\end{abstract}

\tableofcontents
\newpage

% ===========================================================================
\section{为什么必须验证模拟器}

本题的全部结论——完成率、平均定位清除时间、与基线策略的对比——都建立在一个前提上：
\textbf{本地模拟器与官方模拟器在协议和计时上等价}。如果这个前提不成立，
那么论文里所有的数字都只是在描述我自己的程序。

官方模拟器无法获取（百度网盘分发 $+$ 联网注册 $+$ 正式测试额度绑定账号），
因此"验证"必须换一种做法：\textbf{不验证"是否与官方二进制一致"（做不到），
而是验证"是否严格遵守附件 1 与附件 2 的每一条文字规定"}，
并把每一条规定写成一条可执行、可复现、会失败的检查。

\begin{keybox}[title=验证的设计原则]
\textbf{（1）检查必须能失败}：每条检查都断言一个具体数值或状态，而不是"没报错就算过"。\par
\textbf{（2）区分"实现错"与"测试错"}：第一次运行 60 项里失败 10 项，
其中 6 项事后证明是我的测试写错了（把绝对坐标当位移、在已开启的会话上测"未进场"语义等），
只有 4 项是实现缺陷。两类错误必须分开处理，否则会为了"让测试变绿"而改坏实现。\par
\textbf{（3）优化必须自带等价性证明}：任何为了速度而重写的关键逻辑（本文是增量认证扫描），
都要有一条检查与"从零重算"逐次比对。
\end{keybox}

% ===========================================================================
\section{64 项检查的构成}

\begin{center}
\begin{tabular}{clcp{6.6cm}}
\toprule
组 & 覆盖范围 & 项数 & 代表检查 \\ \midrule
A & HTTP 协议 & 18 & 四条路由、\texttt{Content-Type}/\texttt{Encoding} 校验、
请求体上限、JSON 合法性、字段白名单、\texttt{arena\_id}/\texttt{robot\_id} 绑定、
坐标与频道取值域、幂等重放、未进场/重复进场/退出后调用 \\
B & 虚拟时间记账 & 8 & 移动 $=d/5$、切换 $1$ s（同频道 $0$ s）、检测固定 $5$ s、
清除命中 $5$ s/未命中 $3$ s、清除不改频道、时钟单调、已清除源静默 \\
C & 物理判据 & 16 & 源个数 $10\sim16$、频道互异、$R_{\rm eff}\in[1000,1500]$、
面积均匀采样、超距静默、$5$ m 内 \texttt{near}、误差上界、同点重测一致、
清除半径 $20$ m 边界、定向半平面含边界、反侧静默、清除与角度无关 \\
D & 响应与日志契约 & 8 & \texttt{/enter} 回传三个时限、\texttt{measure\_result} 三态、
\texttt{clear\_result} 两态、日志五个必备字段与成本分解 \\
E & 开源参考交叉对照 & 5 & 12 个常量逐项相等、300 步随机动作逐步成本相等、
算例统计量一致、误差模型分歧记录 \\
F & 守恒律 & 4 & 日志分量之和 $=$ 最终虚拟时钟、行程守恒、单局通关、策略时钟不低于竞技场时钟 \\
G & 增量认证等价性 & 5 & 每次扫描与暴力重算逐次比对（60 次）、插桩不改变单局结果、
边界外补测点确实在域外、补测点数受限、内部候选不触发域外 \\
\bottomrule
\end{tabular}
\end{center}

运行方式与产物：
\begin{lstlisting}[language=bash]
cd code
python verify_simulator.py          # 64 项检查，全过返回 0，否则返回 1
# 报告写入 data/simulator_verification.json
\end{lstlisting}

% ===========================================================================
\section{抓到的三个真实缺陷}

\subsection{缺陷 1：\texttt{robot\_id} 绑定形同虚设}

\begin{center}
\begin{tabular}{p{3.2cm}p{11.6cm}}
\toprule
检查 & \texttt{robot\_id\_mismatch}：用错误的 \texttt{robot\_id} 发 \texttt{/measure}，
应当返回 \texttt{200 + accepted=false} \\ \midrule
实测 & \texttt{200 + accepted=true}（被接受了） \\
根因 & HTTP 层只检查\textbf{服务启动时传入}的 \texttt{robot\_id}。
默认启动（不传）时该值为 \texttt{None}，判断被整体跳过 \\
修法 & 竞技场在首次成功 \texttt{/enter} 时绑定该 \texttt{robot\_id}（官方运行时本来就"认识参赛队"），
HTTP 层同时查询启动参数与绑定值 \\
影响 & 不影响本队成绩（我们从不发错 id），但会让模拟器作为"公共 mock"时语义失真 \\
\bottomrule
\end{tabular}
\end{center}

\subsection{缺陷 2：四舍五入突破 $\pm1^{\circ}$ 硬上界}

\begin{center}
\begin{tabular}{p{3.2cm}p{11.6cm}}
\toprule
检查 & \texttt{bearing\_error\_bound}：400 个有效样本的真实误差绝对值必须 $\le1^{\circ}$ \\ \midrule
实测 & 最大 $1.0024^{\circ}$ \\
根因 & 误差在 $\pm1^{\circ}$ 内均匀采样后\textbf{再四舍五入到两位小数}，
边界样本被舍入推到区间外 \\
修法 & 新增 \texttt{\_round\_within\_error()}：舍入后若偏离真值超过上界，
就按 $0.01^{\circ}$ 步长往回收，直到落回界内 \\
影响 & 官方 mock 与我们的实现都有两位小数格式；修正后 20000 个样本最大误差 $0.99985^{\circ}$。
这条修正的意义在于：\textbf{$\pm1^{\circ}$ 是题面的硬约束，不应因为显示格式而被突破} \\
\bottomrule
\end{tabular}
\end{center}

\subsection{缺陷 3：认证扫描的"未认证位置"列表被截断（最有价值的一条）}

认证扫描原本的实现里有一句看似合理的早退：

\begin{lstlisting}[language=Python,caption={原实现中的早退（已移除）}]
for q in grid:
    for c in channels:
        if flags.get(c) is False:      # 该频道已经确定未通过
            continue                   # <-- 跳过它，但 q 也不再被记入 uncov
        ...
        if not ruled_out:
            flags[c] = False
            pending.append(c)
\end{lstlisting}

\textbf{当一个频道在某个候选位置失败后，它在后续所有候选位置都被跳过}，
于是这些位置不会进入 \texttt{uncov}。这意味着：

\begin{itemize}[leftmargin=1.8em]
  \item \texttt{ok = (len(uncov) == 0)} 仍然正确（只要有一个点失败就非空），
        所以\textbf{终止判据没错}；
  \item 但规划器看到的"未认证区域"只是真实区域的很小一部分——
        实测一次调用里真实未认证点 1253 个，而返回的只有 \textbf{1} 个。
\end{itemize}

重写为\textbf{单调增量}版本后（"已认证"只可能从否变是，因此每个候选点一旦被排除就永不回头；
每来一条新读数，只需对仍然开放的候选点重新判定），
既修好了完整性，又把单局 CPU 从 41 s 降到 2.8 s（\textbf{15 倍}）。

\begin{keybox}[title=这条缺陷是怎么被发现的]
它不是我读代码读出来的，而是\textbf{第 G 组检查}逼出来的：
该组要求"增量结果必须等于从零重算的结果"。第一次运行立刻报
\texttt{call 1: brute 1 vs incremental 1253 uncovered}。
追下去才发现\textbf{错的不是新的增量实现，而是我用来当基准的旧逻辑}——
这也说明"写一条与实现独立的参照实现"是抓这类缺陷的唯一可靠办法。
\end{keybox}

% ===========================================================================
\section{与开源实现的交叉对照}

\subsection{调研结果}

按"与本题的可对照程度"排序：

\begin{longtable}{p{3.9cm}p{1.9cm}p{9.0cm}}
\toprule
项目 & 类型 & 与本题的关系 \\ \midrule
\endhead
\texttt{JammersSimulator\hspace{0pt}-\hspace{0pt}Tool}\newline
\scriptsize \texttt{github.com/\hspace{0pt}Jammers\hspace{0pt}-\hspace{0pt}Simulator\hspace{0pt}-\hspace{0pt}Lab} &
MIT，Python & \textbf{直接针对 2026 国赛 B 题}。官方 Windows 模拟器的命令行/MCP 驱动，
并附一个\textbf{离线 mock 模拟器}（\texttt{mock\_simulator.py}）。
README 明确声明"仅含驱动，不含任何解法"，并以双盲评审为由做了脱敏。
它是唯一能拿来逐项对照协议与计时的公开实现 \\
\texttt{dstl/\hspace{0pt}Stone\hspace{0pt}-\hspace{0pt}Soup}\newline \scriptsize UK 国防科技实验室 &
BSD 类，Python & 成熟的跟踪/估计框架，含\textbf{纯方位跟踪}示例
（\texttt{RadarBearing} 传感器 $+$ \texttt{Cartesian2DToBearing} 观测模型 $+$ 扩展卡尔曼滤波）。
对照价值在\textbf{方法论}：它代表"把方位测量喂给贝叶斯滤波器递归估计目标状态"的路线；
本题我们走的是\textbf{确定性集合}路线（角楔求交得到定位区域，最小包围圆给不确定度），
不需要先验与过程模型，且能给出\textbf{可证明的终止判据} \\
\texttt{tiiuae/\hspace{0pt}gnn\hspace{0pt}-\hspace{0pt}jamming\hspace{0pt}-\hspace{0pt}source\hspace{0pt}-\hspace{0pt}localization}\newline
\scriptsize 阿联酋 TII &
研究代码 & 用图神经网络做\textbf{干扰源定位}，属同一应用域（RF 干扰源）但方法为数据驱动，
需要仿真数据集训练；与本题的几何/在线决策路线互补 \\
\texttt{dressel/\hspace{0pt}FEBOL.jl} & Julia 研究代码 &
Filter Exploration for Bearing Only Localization，研究"为了定位该往哪里走"，
和问题二的站位优化是同一类问题 \\
\texttt{El\hspace{0pt}-\hspace{0pt}Mohr/\hspace{0pt}WSN\hspace{0pt}-\hspace{0pt}Loclaization\hspace{0pt}-\hspace{0pt}Simulator} & Python &
无线传感器网络定位仿真环境，可作测距/测角定位的可视化参照 \\
\texttt{daniyalshagiroff/\hspace{0pt}acoustic\hspace{0pt}-\hspace{0pt}triangulation} & Python &
声学三角定位求解器，其\textbf{三角交会求解}与本文问题一的角楔求交是同一数学内核 \\
\bottomrule
\end{longtable}

\begin{warnbox}[title=诚实说明]
除第一项外，其余项目我只核对了其公开说明与示例，\textbf{没有逐行审阅源码}，
因此上面的定位是"可对照性评估"而非"代码级比对结论"。
第一项因为直接对应本题协议，做了逐项常量与成本模型比对（见下）。
\end{warnbox}

\subsection{逐项常量对照：12/12 相等}

\begin{center}
\begin{tabular}{llcc}
\toprule
量 & 含义 & 参考 mock & 本实现 \\ \midrule
速度 & 移动 & $5.0$ m/s & $5.0$ m/s \\
切换 & 频道切换 & $1.0$ s（同频道 $0$） & 同 \\
检测 & 一次测向 & $5.0$ s & 同 \\
清除命中 & 光学 $3$ $+$ 激光 $2$ & $5.0$ s & 同 \\
清除未命中 & 仅光学 & $3.0$ s & 同 \\
\texttt{near} & 过近阈值 & $5.0$ m & 同 \\
清除半径 & 与角度无关 & $20.0$ m & 同 \\
场地半径 & 目标域 & $1800$ m & 同 \\
接收半径 & 每源 & $[1000,1500]$ m & 同 \\
源个数 & 每局 & $[10,16]$ & 同 \\
虚拟时长上限 & --- & $360000$ s & 同 \\
真实时长上限 & --- & $1200$ s & 同 \\
\bottomrule
\end{tabular}
\end{center}

\subsection{成本模型：300 步随机动作逐步相等}

把参考 mock 的计费公式（移动 $d/5$ $+$ 切换 $[0,1]$ $+$ 检测 $5$；
清除 移动 $d/5$ $+$ 光学 $3$ $+$ 激光 $[0,2]$）独立实现为参照，
然后让本模拟器执行 300 次随机 \texttt{/measure} 与约 90 次随机 \texttt{/clear}，
逐步比较虚拟时间增量：

\begin{center}
\begin{tabular}{lc}
\toprule
检查 & 结果 \\ \midrule
\texttt{measure\_cost\_matches\_reference} & 最大偏差 $<10^{-9}$ s \\
\texttt{clear\_cost\_matches\_reference} & 最大偏差 $<10^{-9}$ s \\
\bottomrule
\end{tabular}
\end{center}

两个\textbf{独立实现}在同一串动作上给出逐步相同的虚拟时间——这是"时间记账忠实"最强的证据。

\subsection{唯一分歧：示向度误差模型}

\begin{center}
\begin{tabular}{p{4.0cm}p{5.4cm}p{5.4cm}}
\toprule
 & 参考 mock & 本实现 \\ \midrule
实现 & \texttt{err = rng.uniform(-1,1)} & \texttt{crc32(位置, 频道)} 映射到 $[-1,1]$ \\
每次调用 & \textbf{重新随机} & \textbf{完全相同} \\
附件 1 要求 & \multicolumn{2}{c}{"同一地点重复测量结果相同"} \\
是否符合题设 & \textbf{否} & 是 \\
\bottomrule
\end{tabular}
\end{center}

这一分歧很重要，因为\textbf{它会影响策略选择}：如果每次重测都重新随机，
那么"在同一点多测几次取平均"就能降低误差，站位与信息获取的模型都要改写。
该 mock 的定位是"离线开发与 CI 用的简化桩"（其 README 自述），
我们没有理由跟随它偏离附件 1。

\begin{keybox}[title=本实现的误差模型为何刻意"无空间结构"]
误差是 $(\text{位置},\text{频道})$ 的纯函数，按 $0.01$ m 取整后哈希。
这意味着\textbf{相邻地点的误差完全独立}——现实中测向误差常有空间相关性
（同一建筑反射造成同向偏差）。我们刻意不引入这种结构，
以免结论建立在"利用了现实中并不存在的平滑性"之上。这是一次保守选择。
\end{keybox}

% ===========================================================================
\section{一个把两轮调参全部作废的陷阱}

\subsection{现象}

第一轮坐标下降给出的结果好得可疑：

\begin{lstlisting}[caption={第一轮寻优输出（有问题）}]
[q3] baseline ratio=1.0000 time=324.63
[q3] pass0 ring_radius = 1250.0 -> ratio=1.0000 time=181.06
[q3] FINAL ratio=1.0000 time=181.06
\end{lstlisting}

一个 $1280\to1250$ 的站位半径改动，把平均时间砍掉 44\%？在独立算例池上复算：

\begin{lstlisting}[caption={独立池（种子 31000--31039）复算}]
Q3 baseline (rho=1280)  ratio=1.0000  time= 331.8  travel=14392
Q3 tuned    (rho=1250)  ratio=1.0000  time= 326.2  travel=14097
\end{lstlisting}

\textbf{只有 1.7\%，不是 44\%。}

\subsection{根因与最小复现}

\texttt{sim.Arena(sources)} 保存的是传入的 \texttt{Source} 对象本身（浅拷贝），
运行中被清除的源会把 \texttt{cleared=True} 写在\textbf{共享对象}上。
而寻优脚本为了"每个候选都在同一批算例上评估"，把算例池只生成了一次：

\begin{lstlisting}[language=Python,caption={缺陷写法}]
cases = pool(tag, n_cases)              # 只生成一次
for name, values in grid.items():
    sc = evaluate(cand, cases)          # 每次都在同一批 Source 对象上跑
                                     # -> 第 2 次起，源都已经被清空了
\end{lstlisting}

最小复现（5 个算例，连跑两次）：

\begin{lstlisting}[caption={复用池 vs 重新生成}]
reused pool  eval0: ratio=1.000 time=334.3
reused pool  eval1: ratio=1.000 time=188.2      <-- 假的"提升"
regenerated  eval0: ratio=1.000 time=334.3
regenerated  eval1: ratio=1.000 time=334.3      <-- 真实
\end{lstlisting}

\textbf{第二轮寻优同样中招}。这不是新问题——论文自检报告里记录的故障 1 就是它，
当时修的是"调参脚本"，但新写的寻优脚本又犯了一次。
说明\textbf{光靠"上次修过"不够，必须让脚本自己拒绝这种状态}。

\subsection{修法与守卫}

\begin{enumerate}[leftmargin=1.8em]
  \item \textbf{每次评估都从种子重新生成算例}（生成成本几毫秒，相对单局 1--3 s 可忽略）；
  \item \textbf{加可复现性守卫}：启动时对同一配置连算两次，
        若分数不一致则直接退出（返回码 2）并打印诊断：
\begin{lstlisting}[caption={守卫的输出}]
[guard] ok: repeated evaluation is reproducible (1.0000, 1490.6 s)
\end{lstlisting}
\end{enumerate}

\subsection{修正后的真实收益}

在\textbf{从未参与过调参的独立算例池}上（问题三 种子 12000 起，问题四 种子 62000 起）
比较"优化前配置"与"验证后配置"——即同时计入本轮的三项算法改进
（增量认证、定向场景下的无谓搜索站位抑制、Or-opt）：

\begin{center}
\begin{tabular}{lcccc}
\toprule
 & 完成率 & 最差完成率 & 平均时间/s & 平均行程/m \\ \midrule
问题三 优化前 & 1.0000 & 1.0000 & 见 \S 6 表 & --- \\
问题四 优化前 & 0.9974 & 0.9375 & 1395.6 & 78202 \\
问题四 优化后 & \textbf{1.0000} & \textbf{1.0000} & \textbf{1114.0} & 见 \S 6 \\
\bottomrule
\end{tabular}
\end{center}

结论：\textbf{问题三的性能是结构性的}——8 个参数在 313$\sim$348 s 的窄带内几乎无差异，
说明瓶颈是"覆盖 $+$ 认证"的固有行程，不是参数没调好；
\textbf{问题四的收益是真实的}，而且必须分两步看。

\textbf{第一步：修正的代价。} 修正认证扫描后，"提前宣布完成"不再发生，
机器人必须真正走完搜索，因此时间上升：

\begin{center}
\begin{tabular}{lcccc}
\toprule
 & 完成率 & 最差完成率 & 平均时间/s & 平均行程/m \\ \midrule
修正前（错误认证） & 0.9574 & 0.8333 & 661.0 & 28476 \\
修正后 & \textbf{0.9935} & 0.9091 & 1476.1 & 73064 \\
\bottomrule
\end{tabular}
\end{center}

\textbf{第二步：在正确的前提下重新优化。} 把搜索成本按任务类别拆开后发现：
\textbf{89\% 的时间与行程都花在"搜索站位"上}（65 次搜索任务对 12 次清除任务），
而其中相当一部分站位是"域内永远无法完成认证"的边界候选点。

据此做了两项改动并逐项验收：

\begin{center}
\begin{tabular}{lccc}
\toprule
改动 & 完成率 & 平均时间/s & 平均行程/m \\ \midrule
（基线：修正后） & 0.9944 & 1176.5 & 58944 \\
抑制无法完成认证的搜索站位 & 0.9944 & 955.0 & 47923 \\
再调参（格点间距 1100 m 等） & 0.9935 & \textbf{945} & 47000 \\
\bottomrule
\end{tabular}
\end{center}

\textbf{关键否证}：一度实现的"边界外定向补测"（rim patrol）看似必要
（定理 2 说边界点在域内不可认证），但在\textbf{演练池 30 例}上逐项对照后发现：
开启与关闭的\textbf{完成率完全相同（0.9944，最差也相同 0.9000）}，
而关闭后时间与行程各降 \textbf{18.8\%}。逐个检查失败算例也显示，
漏掉的源半径多在 1158$\sim$1673 m（\textbf{并非贴边}），补测机制对它们毫无帮助。
因此该机制\textbf{默认关闭}，只在需要时作为可选项保留。

\begin{keybox}[title=最终验收（独立算例池，从未参与调参）]
问题四：平均时间 1476.1 s $\to$ \textbf{1018.1 s}（$-31.0\%$），
行程 $-32.8\%$，完成率 0.9935 $\to$ 0.9893（相当）。
问题三：平均时间 333.4 s $\to$ \textbf{323.0 s}（$-3.1\%$），行程 $-9.3\%$，完成率恒为 1.0000。\par
正式测试（HTTP）：问题四由"11/11、11/12、11/11"变为 \textbf{11/11、12/12、11/11}，
此前漏掉的贴边朝外源被成功清除；问题三保持 15/15、12/12、10/10。
\end{keybox}

\textbf{唯一未能消除的失败模式}：贴边且定向方向朝外的源，在域内\textbf{数学上不可认证}
（定理 2）。此时机器人无法证明该频道不存在，只能继续搜索直至预算耗尽。
是否为此付出域外补测的代价（时间约 $\times2.4$）是一个策略选择，
本文默认不付出，并在正文中如实说明。

% ===========================================================================
\section{本轮算法改动清单}

\begin{longtable}{p{3.6cm}p{4.6cm}p{6.6cm}}
\toprule
改动 & 动机 & 验证方式与效果 \\ \midrule
\endhead
增量认证扫描 & 原实现是 $O(\text{网格}\times\text{频道}\times\text{读数})$，
单局 CPU 41 s & 第 G 组逐次与暴力重算比对；
CPU 41 s $\to$ 2.8 s，且修正了 \texttt{uncov} 被截断的缺陷 \\
抑制无谓搜索站位 & 定向场景下，边界候选点附近所有域内方向都指向圆心，
这类站位永远无法完成认证却照样消耗行程 & 10 例 A/B：时间 $-8.5\%$、行程 $-8.7\%$、
完成率不变（1.0000） \\
边界外定向补测 & 边界候选点在域内\textbf{数学上不可认证}（定理 2） & 第 G 组检查补测点确在域外、
数量受限、内部候选不触发；使完成率在困难算例上达到 1.0000 \\
Or-opt 局部搜索 & 最近邻 $+$ 2-opt 之后仍有可改进的短段 &
初期版本在双重循环里全量重算路径长度，导致单局 CPU 数十秒；
改为只优化序列前 12 个任务（滚动时域每次都会重规划，远期任务无意义） \\
\bottomrule
\end{longtable}

% ===========================================================================
\section{如何复现全部验证}

\begin{lstlisting}[language=bash]
cd code
python verify_simulator.py          # 64 项检查 + 报告 data/simulator_verification.json
python optimize_q34.py q3 24 2      # 坐标下降（自带可复现性守卫）
python optimize_q34.py q4 24 2
python final_validate.py 40         # 独立池验收，写 data/final_validation.json
\end{lstlisting}

\begin{warnbox}[title=仍未做到的事]
\textbf{没有官方模拟器上的成绩}。本报告能证明的最强命题是：
"本模拟器在 64 项可检查的协议、计时与物理规定上全部通过，
且其成本模型与一份独立开源实现逐步一致"——这\textbf{不等于}
"在官方模拟器上会得到同样的分数"。若要最终确认，
需要把 \texttt{HTTPClient} 指向官方运行时（默认端口相同），
策略代码无需任何改动。
\end{warnbox}

\end{document}
